%===============================================================
% Chapter 2 — Related Works and Comparative Analysis
% Part A: Sections 2.1–2.10
% Thesis: A Smart Farm IoT Monitoring and Deterministic Agronomic
%         Analytics Platform for Alfalfa Cultivation
% Authors: Abderrahmane Henna · Nacer Eddine Aouissi
% Supervisor: Mohib Eddine KHEBBACHE
% University of El Oued – Echahid Hamma Lakhdar, 2025/2026
%===============================================================

\chapter{Related Works and Comparative Analysis}
\label{chap:related-works}

%---------------------------------------------------------------
\section{Chapter Introduction}
\label{sec:ch2-intro}
%---------------------------------------------------------------

The previous chapter established the theoretical and contextual foundations of
this work: the constraints of Saharan field environments, the data collection
principles that govern a reliable monitoring system, the communication and
topology choices relevant to distributed IoT deployments, and the agronomic
requirements of alfalfa cultivation. Those foundations were drawn from general
literature and technical standards. This chapter moves from general principles to
the specific: it examines how previous research has addressed agricultural IoT
monitoring, what these works have achieved, and where meaningful gaps remain.

The analysis is organized thematically rather than as a sequential list of papers.
Seven dimensions guide the comparison: system architecture and monitoring scope,
wireless communication strategies, soil and environmental data collection,
energy management and field reliability, offline storage and upload handling,
timestamp semantics and data integrity, and dashboard plus analytics design.
A dedicated section then surveys the crop contexts addressed by the reviewed
works, leading to an explicit identification of the alfalfa monitoring gap that
motivates this project.

The chapter draws primarily on three peer-reviewed works that are close enough
to this project to support meaningful comparison: Ahmed et al.
(2022)~\cite{ahmed2022lora_agriculture_chile}, Pereira et al.
(2023)~\cite{pereira2023esp32_drip_irrigation}, and Saban et al.
(2023)~\cite{saban2023smart_agricultural_lora_cloud}. Supporting references
from survey works~\cite{ayaz2019iot_smart_agriculture_survey,farooq2019iot_smart_agriculture_review}
and data quality literature~\cite{teh2020sensor_data_quality,dandrifosse2024weather_qc_agriculture}
provide broader context. Where additional works are referenced and their
bibliographic metadata requires verification, this is noted explicitly.

%---------------------------------------------------------------
\section{IoT-Based Agricultural Monitoring Systems}
\label{sec:iot-monitoring-systems}
%---------------------------------------------------------------

Field-oriented IoT monitoring systems have converged on a recognizable pattern
over the past decade: embedded sensing nodes collect environmental or soil
measurements, a wireless link carries those measurements to a receiver or
gateway, and a web-based interface makes the data accessible to users.
Surveys of the field confirm that this three-layer structure is the dominant
approach~\cite{ayaz2019iot_smart_agriculture_survey,farooq2019iot_smart_agriculture_review}.
Within this shared architecture, individual systems differ significantly in
the variables they monitor, the embedded hardware they use, and the
infrastructure they assume.

ESP32-based nodes appear frequently across recent implementations. Pereira et
al.~\cite{pereira2023esp32_drip_irrigation} present an IoT system built around
the ESP32 microcontroller for smart drip irrigation, integrating soil and
environmental sensing with a cloud-connected dashboard for visualization and
control. The system's primary objective is irrigation automation: the ESP32
measures soil conditions, evaluates them against thresholds, and triggers the
irrigation mechanism accordingly. This is a coherent design for its intended
purpose, and the use of ESP32 as the embedded platform reflects its practical
balance of processing capability, wireless connectivity, and cost. The
limitation relative to this project's goals is one of scope rather than
technical error: the system is designed for irrigation control of a general
crop, not for multi-node field monitoring with multi-domain data governance,
manual agronomic event recording, and an explicit analytical layer targeted at
a specific crop.

Ahmed et al.~\cite{ahmed2022lora_agriculture_chile} approach the monitoring
problem at a different scale. Their LoRa-based platform is deployed across a
large commercial farm in Chile and focuses on remote sensing, LoRa network
behavior, and multi-point data collection. The platform demonstrates that LoRa
communication can sustain reliable data delivery across agricultural field
distances, and the multi-layer system architecture covers sensor nodes, a
gateway, and a server-side monitoring interface. What the work foregrounds is
the feasibility of LoRa as a field communication technology. What it does not
foreground is the data governance layer: the distinction between measurement
time and upload time, the handling of the data record during connectivity
outages, or the reliability of individual readings relative to a quality
control standard.

Saban et al.~\cite{saban2023smart_agricultural_lora_cloud} combine PLC-based
field control with a LoRa/LoRaWAN communication infrastructure and a
cloud-hosted web application for monitoring and device management. The system
extends the monitoring concept toward active agricultural management, and the
cloud web interface provides a relatively complete view of connected devices
and their states. The architecture, however, is fundamentally cloud-first:
field data is expected to reach the server promptly, and the system's
monitoring and control functions operate on this assumption. The treatment of
offline periods, local data persistence, or upload audit trails is not the
central contribution of the work.

Across these systems, the sensing and visualization layers are the primary
technical contributions. The data infrastructure connecting them --- how
measurements are timestamped at the moment of collection, how the local record
is preserved when the network is unavailable, and how uploaded data is
validated against expected structure and timing --- receives comparatively
limited attention. This is not a deficiency in individual papers; it reflects
where the research effort has been directed. It does, however, identify a
space in which the system proposed in this dissertation makes a specific
contribution.

%---------------------------------------------------------------
\section{Wireless Communication in Agricultural IoT}
\label{sec:wireless-communication}
%---------------------------------------------------------------

The appropriateness of LoRa as a field communication technology was established
in Chapter~\ref{chap:background} through a comparison of wireless protocol
characteristics. The reviewed works confirm this choice empirically: both Ahmed
et al.~\cite{ahmed2022lora_agriculture_chile} and Saban et
al.~\cite{saban2023smart_agricultural_lora_cloud} deploy LoRa in agricultural
environments and demonstrate that the technology can sustain data delivery
across field-scale distances with acceptable power budgets.

The practical question that emerges from these works is not whether LoRa can
carry data across a field, but what guarantees accompany delivery. Ahmed et
al.\ demonstrate multi-point remote monitoring with LoRa and discuss network
behavior at the farm scale, including link quality observations. Saban et al.\
use LoRaWAN infrastructure, which adds a network server layer that handles
some aspects of device management and downlink communication. In both cases,
the communication architecture is the subject of design attention. What is less
discussed is the upstream consequence of communication uncertainty: when a
field node fails to transmit successfully, how does the system detect this, log
it, and recover?

Communication feasibility and communication reliability are related but not
identical. A LoRa link that works ninety percent of the time is highly
functional for many monitoring purposes, but the ten percent of missed packets
creates gaps in the server-side record that may not be distinguishable from
periods of genuine field inactivity. If the only signal of a problem is a gap
in the database, an operator looking at a dashboard chart cannot easily
determine whether the crop conditions were stable and unreported, or whether
the node failed and the sensor data was never acquired.

Addressing this requires more than the communication layer alone. An
acknowledgement mechanism at the application level, a diagnostic event log
that records missed transmissions and recovery attempts, and a recovery
procedure for nodes that have drifted out of synchronization with the gateway's
receive window all contribute to turning communication feasibility into
operational reliability. The works reviewed here lay a strong foundation for
the communication design; the proposed system builds on that foundation by
adding the reliability and logging infrastructure on top of the communication
layer.

Regarding outdoor protocol selection more broadly, work on the effects of
communication protocols in precision agriculture outdoor
environments~\cite{iot_protocols_precision_agriculture_outdoor} offers
additional empirical context for understanding how environmental conditions
affect protocol
performance.
% TODO: verify bibliographic metadata for iot_protocols_precision_agriculture_outdoor before final submission

%---------------------------------------------------------------
\section{Soil and Environmental Data Collection Systems}
\label{sec:soil-environmental}
%---------------------------------------------------------------

Soil moisture is the most frequently monitored variable in agricultural IoT
systems, and with good reason: it is the primary determinant of irrigation
timing, a sensitive indicator of crop water status, and a variable that changes
on timescales short enough to make periodic automated measurement practically
useful~\cite{ayaz2019iot_smart_agriculture_survey}. Environmental variables
--- air temperature, humidity, and atmospheric pressure --- appear alongside
soil measurements in many systems as contextual information that helps
interpret the soil data.

The design of the data collection pipeline receives varying degrees of
attention across the reviewed works. Most systems describe the sensing hardware
and the dashboard interface in detail, with the data pathway between them
treated as a solved problem of wireless transmission and database storage.
Wireless sensor network deployments focused specifically on soil moisture
monitoring~\cite{lloret2021soil_moisture_wsn} address aspects of sensor
placement, network topology, and data reliability in that
context.
% TODO: verify bibliographic metadata for lloret2021soil_moisture_wsn before final submission
Studies on practical, low-cost IoT solutions for agricultural data
monitoring~\cite{cheap_practical_iot_agri_monitoring} address deployment
feasibility in resource-constrained
settings.
% TODO: verify bibliographic metadata for cheap_practical_iot_agri_monitoring before final submission

What these systems share is a primary focus on the measurement itself: which
sensor to use, where to place it, and how to transmit its reading. Less
consistently addressed is the organization of the collected data once it
arrives at the server: whether soil and weather data are stored in separate
structured domains, whether individual readings carry quality indicators, and
whether the collection pipeline distinguishes between data that represents a
valid measurement and data that represents a recovered or estimated value.

A related challenge is the multi-source nature of agricultural monitoring
data. Even a relatively simple field setup produces at least two categories of
information: automatic sensor readings and human-entered agronomic records such
as irrigation events and cutting dates. Keeping these categories separate, with
appropriate timestamps for each, allows later analysis to use both without
conflating machine-generated and human-entered information. Teh et
al.~\cite{teh2020sensor_data_quality} identify data completeness and
consistency as critical quality dimensions that must be managed across the full
data lifecycle, not only at the sensor interface. In most reviewed systems,
this lifecycle management is implicit rather than a named design objective.

Pereira et al.~\cite{pereira2023esp32_drip_irrigation} combine soil sensing
with irrigation control in a system that uses sensor measurements to drive
direct actuation. The monitoring and control functions are tightly coupled,
which is appropriate for an irrigation automation system. For a system focused
on monitoring and decision support for a specific crop, however, the separation
between what the system measures and what the operator decides remains
important. Coupling measurement to automatic action blurs this distinction.

%---------------------------------------------------------------
\section{Energy Management and Field Reliability}
\label{sec:energy-reliability}
%---------------------------------------------------------------

Battery-powered field nodes are a practical reality in distributed agricultural
monitoring, and managing power consumption is a standard design challenge. The
reviewed works acknowledge this constraint at different levels of depth.

At the hardware selection level, the choice of microcontroller, radio module,
and sensor type all affect the power budget. ESP32-based systems, as used by
Pereira et al.~\cite{pereira2023esp32_drip_irrigation}, benefit from the
availability of deep sleep modes that reduce current draw significantly during
idle periods. LoRa radios, as used by both Ahmed et
al.~\cite{ahmed2022lora_agriculture_chile} and Saban et
al.~\cite{saban2023smart_agricultural_lora_cloud}, are selected partly for
their low transmission power relative to cellular alternatives.

Energy-aware firmware design is less consistently detailed in the reviewed
works. Duty cycling --- the pattern of wake, measure, transmit, sleep --- is
implied by most field node designs but is not always described in terms of its
effect on measurement continuity or its interaction with the communication
window. The relationship between the node's sleep interval, the gateway's
receive window, and the risk of timing drift between the two is a specific
firmware engineering concern that is rarely foregrounded as a primary design
topic.

A related aspect of field reliability is the system's ability to recover from
communication failures without manual intervention. Pereira et
al.~\cite{pereira2023esp32_drip_irrigation} address retry mechanisms at the
sensor-actuator level. Ahmed et al.~\cite{ahmed2022lora_agriculture_chile}
discuss network performance and link quality. Neither work centers on the
problem of a node that has lost synchronization with the gateway's receive
window and must execute a defined recovery procedure before normal operation
can resume. This distinction between temporary packet loss and timing
desynchronization is subtle but consequential in systems that depend on
coordinated duty cycles.

A broader review of data collection architectures in IoT
networks~\cite{abbaari2024data_collection_iot_networks} discusses reliability
challenges at the network and protocol level, confirming that reliability is a
recognized concern in the IoT literature, even when specific firmware-level
recovery procedures are not detailed in individual
systems.
% TODO: verify bibliographic metadata for abbaari2024data_collection_iot_networks before final submission
A review of smart sensors and their deployment challenges in agricultural
contexts~\cite{rajak2023iot_smart_sensors_agriculture} similarly identifies
reliability and energy as ongoing concerns without providing field-level
recovery
procedures.
% TODO: verify bibliographic metadata for rajak2023iot_smart_sensors_agriculture before final submission

Diagnostic event logging complements energy and reliability design. When a
node misses its expected window, when a recovery attempt is made, or when a
communication failure is resolved, recording these events creates an auditable
history that can inform both real-time dashboard diagnostics and retrospective
reliability analysis. In most reviewed systems, such logging is either absent
or described as a debugging tool rather than a primary feature of the data
model.

%---------------------------------------------------------------
\section{Offline Storage and Upload Strategies}
\label{sec:offline-storage}
%---------------------------------------------------------------

Many agricultural IoT systems are designed with the implicit assumption that
a network connection is available for data forwarding on a timescale
comparable to the measurement interval. Cloud-first architectures, in which
sensor readings are transmitted to a remote server shortly after collection,
are practical when connectivity is reliable. In Saharan and remote field
environments of the kind described in Section~\ref{sec:saharan-agriculture},
this assumption is not always valid.

Ahmed et al.~\cite{ahmed2022lora_agriculture_chile} deploy their platform
across a large commercial farm with LoRa gateways and server-side
infrastructure in place. The architecture assumes that the gateway has
connectivity to forward data to the cloud server. Saban et
al.~\cite{saban2023smart_agricultural_lora_cloud} use a LoRaWAN network server
and a cloud-hosted web application, again assuming that the gateway is
continuously connected to the upstream service. Pereira et
al.~\cite{pereira2023esp32_drip_irrigation} connect the ESP32 to a cloud
service for dashboard access. In each case, the data path leads promptly from
field to cloud.

The consequence of this design is that the server-side record is only as
complete as the connectivity allows. A period of network unavailability
produces a gap in the database that may be indistinguishable from a period
where no field data existed. An offline-first design inverts this priority:
measurements are stored locally before any upload attempt, and the upload
is an explicit, operator-controlled action rather than an automatic background
process. Under this model, the local record is always complete; the server-side
record is a subset of the local record determined by when uploads occurred.

The distinction between when data is uploaded and when it was measured is
therefore not merely a technical detail. It is the design choice that
determines whether a gap in the server database represents an absence of
field data or a delay in its transfer. Treated differently in analytics, these
two interpretations lead to different conclusions about field conditions and
system reliability.

%---------------------------------------------------------------
\section{Time Synchronization and Data Integrity}
\label{sec:time-sync}
%---------------------------------------------------------------

Every data record in a time-series monitoring system carries a timestamp, and
the meaning of that timestamp determines what the record represents. Whether
the timestamp records the moment of measurement, the moment of local storage,
the moment of transmission, or the moment of server-side ingestion is not
always explicitly defined in the reviewed works. In systems with low latency
between measurement and cloud arrival, the difference is small enough to be
negligible. In systems where data can rest on local storage for hours or days
before upload, the distinction becomes analytically significant.

Dandrifosse et al.~\cite{dandrifosse2024weather_qc_agriculture} illustrate
this in the context of agricultural weather data: quality control algorithms
that rely on observation timestamps produce incorrect results when those
timestamps reflect processing time rather than measurement time. The point
applies equally to soil sensor data. A soil moisture reading labelled with the
upload timestamp rather than the measurement timestamp will appear on a chart
at the wrong point in time, potentially misrepresenting the relationship
between soil conditions and agronomic events such as irrigation.

The reviewed works address timestamp handling with varying degrees of
explicitness. Ahmed et al.~\cite{ahmed2022lora_agriculture_chile} focus on the
communication layer and the monitoring platform; the question of which
timestamp propagates through the system and serves as the chart time axis is
not a primary discussion topic. Saban et
al.~\cite{saban2023smart_agricultural_lora_cloud} focus on control and
monitoring through the cloud interface; timestamp semantics at the field node
level are similarly not foregrounded. Pereira et
al.~\cite{pereira2023esp32_drip_irrigation} focus on the irrigation control
logic; the time axis of the sensor readings is not the central concern.

In a system where offline storage and batch upload are by design, timestamp
integrity requires an active approach. The field node must assign measurement
timestamps at the moment of sensing, using a time source that does not depend
on network availability. A dedicated real-time clock module, synchronized
periodically to a reference source and used consistently as the timestamp for
every measurement, provides this guarantee. The clock synchronization
procedure itself must be explicit and auditable, so that an operator can
confirm when the clock was last synchronized and what the expected drift is
between synchronization events.

This approach --- authoritative field-assigned measurement timestamps,
preserved through local storage and batch upload, and used as the sole time
axis for analytics --- is the design choice this dissertation makes and that
the reviewed works do not explicitly address.

%---------------------------------------------------------------
\section{Dashboards, Analytics, and AI Interpretation}
\label{sec:dashboards-analytics-ai}
%---------------------------------------------------------------

Web-based dashboard interfaces appear in all three primary reference systems
as the final layer of user interaction. Ahmed et
al.~\cite{ahmed2022lora_agriculture_chile} provide a monitoring platform with
sensor data visualization. Saban et
al.~\cite{saban2023smart_agricultural_lora_cloud} include a cloud-hosted web
interface for monitoring and device control. Pereira et
al.~\cite{pereira2023esp32_drip_irrigation} implement a dashboard for
visualizing sensor readings and managing irrigation parameters.

These interfaces fulfill the primary function of a monitoring dashboard:
making field data visible to the operator through charts, status indicators,
and historical views. Where they differ from the dashboard proposed in this
work is in the depth of the analytics layer between raw data and displayed
information. A typical monitoring dashboard passes measurements from the
database to the chart directly, possibly with basic filtering and
aggregation. A dashboard backed by a deterministic analytics layer first
applies quality control rules to identify anomalous or unreliable readings,
computes reliability indicators based on the completeness and consistency of
the data record, and makes the results of both the raw and the processed data
visible to the user.

The distinction between raw measurement display and processed data
transparency matters for two reasons. First, an operator looking at a soil
moisture chart cannot easily distinguish a genuine measurement from one that
is the result of a sensor glitch, a duplicate upload, or a timestamp anomaly
--- unless the dashboard makes this distinction explicit. Second, when
analytics outputs are derived from cleaned and quality-controlled data, the
chain from raw reading to displayed result should be traceable, so that a
discrepancy between a sensor reading and an analytics indicator can be
investigated rather than simply accepted.

Regarding artificial intelligence: some recent agricultural systems have
explored machine learning components for prediction, anomaly detection, or
crop state estimation. This is a legitimate research direction, but it
introduces an important boundary question. When AI generates a result ---
whether a prediction, a recommendation, or a summary --- the basis for that
result should be distinguishable from the basis for deterministic metrics
computed by rule-based or statistical methods. Treating AI output as
equivalent to a sensor reading or a reliability score conflates two
fundamentally different types of information.

The approach adopted in this project positions the AI component strictly as
an interpretation layer. Deterministic analytics --- quality control,
reliability scoring, and trend metrics --- are computed independently of any
AI model. The AI component receives the outputs of these deterministic steps
and provides a narrative explanation. It does not compute primary metrics,
fill missing measurements, or override reliability assessments. This boundary
is intentional and is reflected throughout the system architecture.

%---------------------------------------------------------------
\section{Previous Crop Targets in Existing Works}
\label{sec:crop-targets}
%---------------------------------------------------------------

Reviewing the agricultural context of the selected related works reveals a
consistent pattern. Most systems target general irrigation applications,
unspecified crop environments, or broad smart agriculture goals. None of the
primary works reviewed for this chapter targets alfalfa as the principal
monitored crop. Table~\ref{tab:related-work-crop-focus} summarizes the crop
focus and main technical contribution of each reviewed work.

\begin{table}[H]
  \centering
  \caption{Crop context and primary technical focus of selected related works
           reviewed in this chapter.}
  \label{tab:related-work-crop-focus}
  \renewcommand{\arraystretch}{1.3}
  \begin{tabular}{|p{3.8cm}|p{2.4cm}|p{3.6cm}|p{1.4cm}|p{2.0cm}|}
    \hline
    \textbf{Work} & \textbf{Crop / context} & \textbf{Main technical focus} &
    \textbf{Alfalfa target?} & \textbf{Thesis relevance} \\
    \hline
    Ahmed et al.~\cite{ahmed2022lora_agriculture_chile} (2022) &
    Large-scale farms, Chile &
    LoRa IoT monitoring platform &
    No & LoRa architecture comparison \\
    \hline
    Pereira et al.~\cite{pereira2023esp32_drip_irrigation} (2023) &
    General irrigation &
    ESP32 smart drip irrigation &
    No & Embedded sensing comparison \\
    \hline
    Saban et al.~\cite{saban2023smart_agricultural_lora_cloud} (2023) &
    General smart agriculture &
    PLC + LoRa/LoRaWAN + cloud web app &
    No & Cloud dashboard comparison \\
    \hline
    Ayaz et al.~\cite{ayaz2019iot_smart_agriculture_survey} (2019) &
    Multiple crops (survey) &
    IoT smart agriculture survey &
    No & Smart agriculture background \\
    \hline
    Farooq et al.~\cite{farooq2019iot_smart_agriculture_review} (2019) &
    Multiple crops (review) &
    IoT smart farming review &
    No & Smart farming context \\
    \hline
    Rajak et al.~\cite{rajak2023iot_smart_sensors_agriculture} &
    % TODO: verify bibliographic metadata for rajak2023iot_smart_sensors_agriculture before final submission
    Generic IoT agriculture &
    Smart sensors review &
    No & Sensor deployment context \\
    \hline
    Villa-Henriksen et al.~\cite{villa2020iot_arable_farming} &
    % TODO: verify bibliographic metadata for villa2020iot_arable_farming before final submission
    Arable farming (general) &
    IoT challenges and opportunities &
    No & Broad IoT farming review \\
    \hline
    Abbaari et al.~\cite{abbaari2024data_collection_iot_networks} &
    % TODO: verify bibliographic metadata for abbaari2024data_collection_iot_networks before final submission
    IoT networks (general) &
    Data collection architecture &
    No & Data collection theory \\
    \hline
    Cheap practical IoT~\cite{cheap_practical_iot_agri_monitoring} &
    % TODO: verify bibliographic metadata for cheap_practical_iot_agri_monitoring before final submission
    Generic agricultural &
    Low-cost IoT monitoring &
    No & Practical IoT context \\
    \hline
    Lloret et al.~\cite{lloret2021soil_moisture_wsn} &
    % TODO: verify bibliographic metadata for lloret2021soil_moisture_wsn before final submission
    Soil moisture monitoring &
    Wireless soil moisture WSN &
    No & Soil sensing comparison \\
    \hline
  \end{tabular}
\end{table}

The observation that none of the reviewed works targets alfalfa is not in
itself a gap; most monitoring systems are intentionally general and their
value does not depend on crop specificity. The gap emerges when considered
alongside the monitoring requirements established in
Section~\ref{subsec:alfalfa-monitoring-summary}: the repeated cutting cycle,
the soil moisture dynamics between cuts, the need for contextual alignment
between sensor trends and agronomic events, and the specific field conditions
of Saharan oasis cultivation.

%---------------------------------------------------------------
\section{The Alfalfa Monitoring Gap}
\label{sec:alfalfa-gap}
%---------------------------------------------------------------

The practical significance of a monitoring system depends partly on whether
its design reflects the specific requirements of the crop it serves. A system
built for general irrigation scheduling imposes no particular structure on how
field data should be organized around cutting events, how soil moisture trends
before and after an irrigation event should be captured and compared across
different irrigated zones, or how the monitoring timeline aligns with the
biological cycle of the crop. For alfalfa, these are not abstract concerns ---
they are the questions that a monitoring platform must help an operator answer.

Alfalfa's repeated harvest pattern means that the monitoring cycle is not
seasonal in the conventional sense. Multiple cut events occur within a single
growing season, each resetting the soil moisture dynamic and creating a new
period of vegetative regrowth that must be monitored for irrigation timing.
Chapter~\ref{chap:background} established that alfalfa is sensitive to soil
moisture deficit, heat stress, and soil electrical conductivity changes, and
that each of these variables can be monitored with appropriate field sensors.
What the background chapter could not establish from the literature is whether
any existing IoT monitoring system has been designed around this crop in a
Saharan or Algerian context.

Within the works reviewed for this thesis, no selected related work
explicitly targets alfalfa as the monitored crop. The works reviewed focus on
general irrigation systems, unspecified commercial crops, or broad smart
agriculture goals in geographic and climatic settings different from the
Saharan oasis environment. Some work on IoT enhancements specific to Algerian
desert agriculture~\cite{iot_enhancements_algerian_desert_agriculture} touches
on the regional context,
% TODO: verify bibliographic metadata for iot_enhancements_algerian_desert_agriculture before final submission
but does not address alfalfa monitoring as a specific design target. This
observation is stated with care: it reflects the scope of this literature
review, not a claim about all published work globally.

The gap is not only agronomic. Several technical characteristics of the
proposed system respond to limitations identified across the reviewed works:
the assumption of continuous connectivity, the limited treatment of offline
storage and batch upload, the absence of explicit timestamp semantics
distinguishing measurement from transfer, the lack of a diagnostics layer for
recovery and reliability, and the absence of a structured analytics pipeline
with AI bounded to interpretation. Together, these gaps define the space in
which the proposed system makes its contribution.

The combination --- alfalfa as the target crop, Saharan field conditions as
the deployment context, offline-first data collection, RTC-governed measurement
timestamps, deterministic analytics, and a transparent processed-data layer ---
has not been addressed in the reviewed works. The following sections of this
chapter develop the comparative analysis that positions the proposed system
relative to these identified gaps, and the subsequent chapter presents its
design and implementation in detail.

%---------------------------------------------------------------
\section{Proposed System Improvements}
\label{sec:proposed-improvements}
%---------------------------------------------------------------

The gaps identified across the preceding sections --- connectivity assumptions,
limited local persistence, undefined timestamp semantics, absent recovery
procedures, simple analytics, undefined AI boundaries, and no alfalfa-specific
design target --- collectively map the design space that the system proposed in
this dissertation occupies. Each gap corresponds to a deliberate design
decision. This section summarizes those decisions at the level appropriate to
a comparative analysis; their implementation is presented in detail in
Chapter~\ref{chap:implementation}.

The field architecture of the proposed system is organized around three nodes
connected through a LoRa node-to-gateway topology. A MAIN gateway node
coordinates communication, local storage, and upload operations. A dedicated
soil monitoring node (Node~2) acquires soil moisture, soil temperature, and
soil electrical conductivity independently from each irrigated zone. A
dedicated weather node (Node~3) acquires air temperature, relative humidity,
and atmospheric pressure, providing environmental context shared across the
monitored areas. This three-node design, with the gateway coordinating rather
than all nodes communicating as peers, reflects both the field geometry of the
deployment site and the need to configure each node's power budget and duty
cycle independently. A LoRa link with an application-level acknowledgement
mechanism connects the remote nodes to the MAIN gateway; a node that does not
receive an acknowledgement within its expected window enters a defined recovery
procedure rather than failing silently.

Offline-first data acquisition addresses the connectivity constraint that
cloud-first systems cannot resolve by design. All sensor readings are stored
locally on an SD card, in structured CSV format, before any upload is
initiated. System diagnostic events are stored in parallel as JSONL records.
The upload pathway is operator-controlled: a long continuous push-button press
enters Upload Mode, initiating a batch transfer of locally stored records to
the backend server. This design choice carries a direct consequence for time
semantics: the upload timestamp records when the transfer occurred, not when
the measurements were made. Upload time is treated throughout the system as
transfer and audit metadata. It is recorded and available for operational
diagnosis, but it does not function as the time axis for any sensor chart or
analytics output.

The temporal integrity of measurement timestamps requires a reliable time
source that is independent of network availability. A DS3231 real-time clock
module, integrated into the MAIN gateway node, provides this source. The
DS3231 maintains time across power cycles without requiring a network
connection. Clock synchronization to an external reference is performed
explicitly: pressing the field push-button four consecutive times enters RTC
Sync Mode, during which the clock is updated and the synchronization event is
logged. The \texttt{measured\_at} field derived from this clock is the
canonical timestamp for every sensor record stored on SD, transferred in batch
uploads, and used in the database for all chart time axes, analytics
computations, and alert threshold evaluations.

Firmware operation is structured around four explicit modes in addition to the
normal measurement cycle: Upload Mode for batch data transfer, Recovery Mode
for nodes that have lost synchronization with the gateway's receive window, RTC
Sync Mode for clock maintenance, and Power Saving Mode for battery-constrained
periods. Power Saving Mode employs configured deep sleep intervals between
measurement cycles and uses relay-controlled switching to cut power to the
soil sensor between readings, reducing quiescent current draw. Every mode
transition, communication failure, recovery attempt, and synchronization event
is written to the JSONL diagnostic event log on the SD card. This log is
uploaded alongside sensor readings and made available to the dashboard
diagnostics interface, producing an auditable history of system behavior that
operators can examine alongside sensor trends.

The analytics layer is structured as a deterministic pipeline executed on the
backend after each upload is processed. Quality control rules identify
anomalous readings, structural inconsistencies, and timestamp irregularities
within the uploaded batch. A cleaning step separates accepted records from
flagged ones without modifying the raw data. Reliability scoring evaluates the
completeness and internal consistency of the data record within each analysis
window. Alert rules flag readings that cross configured threshold values.
Analytics snapshots consolidate these outputs into structured representations
consumed by the dashboard. A dedicated Processed Data Viewer exposes the full
traceability chain: the raw reading, the quality-control result, and the
processed representation are all accessible to the operator. No raw data is
mutated at any stage of this pipeline, so any discrepancy between a sensor
reading and an analytics output can be investigated and explained.

The AI component of the system operates as a bounded interpretation layer.
The Gemini model is accessed exclusively through a backend proxy and is never
given raw sensor data directly. It does not execute database queries, compute
reliability scores, or evaluate quality control rules. The input to the AI
component is a structured snapshot prepared from the deterministic analytics
outputs, applied after quality control and reliability assessment have been
completed. From this snapshot, the AI generates a narrative interpretation.
Deterministic metrics, quality control flags, and reliability scores remain the
authoritative outputs; AI-generated text is supplementary and is presented
distinctly in the dashboard interface. This boundary is an architectural
decision, not a limitation, and it ensures that the interpretive layer cannot
override or substitute for the computed layer. Agronomic decisions remain with the human operator; AI-generated narratives function as decision-support information, not as agronomic prescriptions.

Crop specificity is embedded in the monitoring design at the level of sensor
variable selection, agronomic event structure, and analytics parameterization.
Alfalfa (\textit{Medicago sativa}) is the target crop for the deployment field
in the El Oued region of Algeria. The monitoring requirements established in
Chapter~\ref{chap:background} --- cut-cycle tracking, soil moisture dynamics
between irrigation events, soil electrical conductivity as a relative trend
indicator, and sensitivity to heat and water deficit during regrowth --- guided
the selection of monitored variables and the design of the agronomic event log.
Within the works selected for review in this chapter, no comparable system
addresses these requirements for alfalfa in a Saharan field environment. The
proposed system provides a field-oriented prototype that responds to this
combined crop and technical gap.

%---------------------------------------------------------------
\section{Comparative Analysis}
\label{sec:comparative-analysis}
%---------------------------------------------------------------

The comparison between the selected related works and the proposed system is
developed across three complementary views: a grouped crop-focus summary that
extends the survey in Section~\ref{sec:crop-targets} by positioning the
proposed system alongside the reviewed works, a criterion-based table that
examines twelve specific design dimensions where the proposed system diverges
from prior approaches, and a binary feature matrix that enables direct
comparison across fifteen system-level features.

\subsection*{Crop Focus and Agricultural Context}

Table~\ref{tab:crop-focus-summary} extends the per-work crop survey of
Section~\ref{sec:crop-targets} by grouping the reviewed works into categories
and adding the proposed system as a final row, making explicit the aggregate
gap: within the selected works reviewed for this thesis, no prior system
explicitly targets alfalfa as the main monitored crop.

% NOTE: This table uses label tab:crop-focus-summary to avoid duplication
% with tab:related-work-crop-focus (already defined in Section 2.9).
% Both tables address crop context; this one groups works and adds the
% proposed system row for comparative positioning.
\begin{table}[H]
  \centering
  \caption{Grouped crop focus and agricultural context of reviewed works compared
           with the proposed system.}
  \label{tab:crop-focus-summary}
  \renewcommand{\arraystretch}{1.3}
  \begin{tabular}{|p{3.5cm}|p{2.8cm}|p{2.8cm}|p{1.5cm}|p{2.0cm}|}
    \hline
    \textbf{Work} & \textbf{Agricultural context} &
    \textbf{Crop focus} & \textbf{Alfalfa target?} & \textbf{Thesis relevance} \\
    \hline
    Ahmed et al.~\cite{ahmed2022lora_agriculture_chile} (2022) &
    Large commercial farms, Chile &
    General farm monitoring &
    No & LoRa architecture comparison \\
    \hline
    Pereira et al.~\cite{pereira2023esp32_drip_irrigation} (2023) &
    Irrigated field, unspecified crop &
    Drip irrigation application &
    No & Embedded sensing comparison \\
    \hline
    Saban et al.~\cite{saban2023smart_agricultural_lora_cloud} (2023) &
    General smart agriculture &
    Unspecified crop context &
    No & Cloud monitoring comparison \\
    \hline
    Survey works~\cite{ayaz2019iot_smart_agriculture_survey,farooq2019iot_smart_agriculture_review} &
    Multiple regions and crops &
    Broad IoT agriculture &
    No & Background and framing \\
    \hline
    Other reviewed works [TODO] &
    Various regional or general contexts &
    Generic or unspecified crops &
    No & Contextual support \\
    \hline
    \textbf{Proposed system} &
    Saharan field, El Oued, Algeria &
    Alfalfa (\textit{Medicago sativa}) &
    \textbf{Yes} & --- \\
    \hline
  \end{tabular}
\end{table}

Within the reviewed works used in this thesis, no selected work explicitly
targets alfalfa as the main monitored crop. The proposed system is distinctive
in this respect: alfalfa cultivation defines not only the deployment crop but
also the sensor variable selection, the structure of the agronomic event log,
and the analytics parameterization applied to the field data.

\subsection*{Criterion Comparison}

Table~\ref{tab:related-work-comparison} examines the reviewed works against the
proposed system across twelve design dimensions. Each dimension corresponds to
a thematic gap identified in Sections~\ref{sec:iot-monitoring-systems}
through~\ref{sec:dashboards-analytics-ai}.

\begin{table}[H]
  \centering
  \caption{Criterion-based comparison of selected related works and the proposed
           system across twelve design dimensions.}
  \label{tab:related-work-comparison}
  \renewcommand{\arraystretch}{1.3}
  \begin{tabular}{|p{2.5cm}|p{3.0cm}|p{3.3cm}|p{3.3cm}|}
    \hline
    \textbf{Criterion} & \textbf{Selected related works} &
    \textbf{Observed limitation} & \textbf{Proposed system response} \\
    \hline
    Target crop &
    General irrigation, mixed or unspecified crops &
    No alfalfa-specific design, event model, or agronomic context &
    Alfalfa target; cut-cycle-aware agronomic event log \\
    \hline
    Deployment context &
    Commercial farms, general agricultural settings &
    Limited attention to Saharan constraints: unstable power, limited connectivity &
    Field-oriented prototype, El Oued oasis environment \\
    \hline
    Communication &
    LoRa/LoRaWAN (Ahmed, Saban); cloud Wi-Fi path (Pereira) &
    Feasibility demonstrated; application-level recovery not foregrounded &
    LoRa node-to-gateway with ACK; defined Recovery Mode for desynchronization \\
    \hline
    Offline operation &
    Cloud-first or continuously connected architectures assumed &
    Database gaps indistinguishable from absent field data &
    SD-first storage; server record is an explicit subset of local record \\
    \hline
    Upload strategy &
    Automated or continuous cloud upload &
    Upload time conflated with measurement time in downstream analysis &
    Manual batch upload (long button press); upload time = transfer metadata only \\
    \hline
    Time policy &
    Timestamp semantics not explicitly defined &
    Risk of charting data against processing or arrival time &
    DS3231 RTC; \texttt{measured\_at} = canonical time axis for all analytics \\
    \hline
    Power management &
    Hardware selected for low power; firmware duty cycle often implicit &
    Desynchronization consequence of duty-cycle drift not addressed &
    Power Saving Mode: deep sleep intervals + relay-controlled sensor power \\
    \hline
    Recovery &
    Link-layer packet retries; no application-level recovery procedure &
    Missed windows accumulate; node timing state not explicitly recoverable &
    Recovery Mode: defined firmware procedure for timing desynchronization \\
    \hline
    Data domains &
    Single or undifferentiated data stores &
    Sensor readings, events, and uploads may be mixed or conflated &
    Four-domain schema: sensor readings, system events, uploads, agronomic events \\
    \hline
    Dashboard scope &
    Sensor reading visualization; device status in some systems &
    Raw chart display; processing steps not visible to operator &
    Processed Data Viewer: raw, cleaned, and processed lineage accessible \\
    \hline
    Analytics &
    Basic aggregation or threshold alerts in some systems &
    No explicit quality control, reliability scoring, or data-quality visibility &
    Deterministic QC, reliability scoring, alert rules, analytics snapshots \\
    \hline
    AI role &
    Not addressed, or general AI potential noted &
    Risk of conflating AI output with sensor or metric authority &
    Gemini via backend proxy; interpretation only; deterministic metrics authoritative \\
    \hline
  \end{tabular}
\end{table}

\subsection*{Feature Matrix}

Table~\ref{tab:feature-matrix} provides a feature-level summary comparing the
three primary reviewed works and the proposed system across fifteen
system-level capabilities. Survey and unverified works are grouped as
``Other reviewed'' and assessed where feature evidence is clearly available.

\begin{table}[H]
  \centering
  \caption{System feature comparison across selected related works and the
           proposed system. \checkmark~= present; \texttt{--}~= absent;
           Partial~= partially addressed; Not specified~= feature not discussed
           in the published work; [TODO]~= metadata unverified.}
  \label{tab:feature-matrix}
  \resizebox{\textwidth}{!}{%
  \renewcommand{\arraystretch}{1.35}
  \begin{tabular}{|p{5.2cm}|c|c|c|c|c|}
    \hline
    \textbf{Feature} &
    \textbf{Ahmed et al.} &
    \textbf{Pereira et al.} &
    \textbf{Saban et al.} &
    \textbf{Other reviewed} &
    \textbf{Proposed} \\
    \hline
    LoRa communication &
    \checkmark & -- & \checkmark & Partial & \checkmark \\
    \hline
    Multi-node field architecture &
    \checkmark & -- & \checkmark & Partial & \checkmark \\
    \hline
    Soil sensing &
    \checkmark & \checkmark & \checkmark & \checkmark & \checkmark \\
    \hline
    Weather sensing &
    \checkmark & Partial & \checkmark & Partial & \checkmark \\
    \hline
    SD local storage &
    -- & -- & -- & -- & \checkmark \\
    \hline
    Manual batch upload &
    -- & -- & -- & -- & \checkmark \\
    \hline
    RTC-based measurement timestamp &
    -- & -- & -- & Not specified & \checkmark \\
    \hline
    Firmware Recovery Mode &
    -- & -- & -- & -- & \checkmark \\
    \hline
    Power Saving Mode (deep sleep + relay) &
    -- & Not specified & -- & -- & \checkmark \\
    \hline
    Diagnostic event logging &
    -- & -- & -- & -- & \checkmark \\
    \hline
    Deterministic QC and data cleaning &
    -- & -- & -- & -- & \checkmark \\
    \hline
    Reliability scoring &
    -- & -- & -- & -- & \checkmark \\
    \hline
    Processed Data Viewer (lineage) &
    -- & -- & -- & -- & \checkmark \\
    \hline
    AI interpretation boundary &
    -- & -- & -- & -- & \checkmark \\
    \hline
    Alfalfa-specific monitoring &
    -- & -- & -- & -- & \checkmark \\
    \hline
  \end{tabular}%
  }
\end{table}

The matrix illustrates that the primary reviewed works demonstrate strong
capability in field communication infrastructure, embedded sensing, and
dashboard visualization, dimensions on which the proposed system also builds.
The proposed system's extension lies in the data governance and analytics layers
that sit between sensor acquisition and operator display: SD-based offline
storage and manual batch upload as a connectivity-resilient pipeline;
RTC-authoritative measurement timestamps preserved through local storage and
batch transfer; a defined firmware procedure for recovery from timing
desynchronization; diagnostic event logging as a first-class data domain;
a deterministic analytics pipeline with quality control, reliability scoring,
and processed-data transparency; a constrained AI interpretation layer; and
alfalfa as an explicit crop target with Saharan deployment context. These are
absent from all three primary comparisons and, to the extent observable, from
the broader reviewed literature.

%---------------------------------------------------------------
\section{Chapter Summary}
\label{sec:ch2-summary}
%---------------------------------------------------------------

This chapter has examined published agricultural IoT monitoring systems
through seven thematic dimensions: monitoring architecture and scope, wireless
communication, soil and environmental data collection, energy management and
field reliability, offline storage and upload strategies, time synchronization
and data integrity, and dashboard with analytics design. Three peer-reviewed
works provided the primary comparison
basis~\cite{ahmed2022lora_agriculture_chile,pereira2023esp32_drip_irrigation,saban2023smart_agricultural_lora_cloud},
supported by survey
literature~\cite{ayaz2019iot_smart_agriculture_survey,farooq2019iot_smart_agriculture_review}
and data quality
references~\cite{teh2020sensor_data_quality,dandrifosse2024weather_qc_agriculture}.
Additional works referenced in the thematic sections carry pending metadata
verification and are not cited as primary evidence.

The reviewed works represent substantive contributions to agricultural IoT
research. LoRa has been demonstrated as a suitable communication medium for
large-scale field monitoring. ESP32-based embedded nodes have been deployed
effectively for irrigation-oriented sensing. Cloud-connected web applications
give agricultural operators remote visibility over their field installations.
These results are established, and the proposed system builds on them rather
than contesting them.

The thematic analysis reveals that the layers between sensor acquisition and
operator display receive comparatively less attention in the reviewed works.
Local persistence strategies that sustain data collection through connectivity
outages are not the central design objective of the primary works, which assume
prompt or near-prompt cloud delivery. Timestamp discipline --- the practice of
assigning measurement time at the sensor and carrying that timestamp
unchanged through storage, upload, and analytics --- is not explicitly
addressed in any of the primary works. Recovery procedures that return a field
node to synchronized operation after a timing desynchronization event, as
opposed to link-layer packet retries, are absent. Quality control and
reliability scoring are not presented as distinct pipeline stages, and the
chain from raw reading to displayed result is not exposed to the operator as
a traceable artefact. The AI boundary question, concerning what an AI component
may and may not do within a data-driven monitoring system, is not addressed in
the primary works. Alongside these technical dimensions, no reviewed work
targets alfalfa cultivation in a Saharan field environment, leaving both a crop
and a deployment context unaddressed.

The proposed system is not positioned as a universal improvement over all prior
work. Its contribution is specific: stronger data governance through
offline-first acquisition, authoritative field-assigned measurement timestamps,
domain-separated storage, a deterministic quality control and reliability
pipeline, processed-data traceability, and a bounded AI interpretation layer,
applied to alfalfa cultivation in the El Oued region of Algeria. Chapter~\ref{chap:implementation}
presents the design and implementation of the system components that realize
these contributions, from the distributed field hardware and firmware
operational modes through the backend analytics pipeline and the dashboard
interface that makes the results accessible to the operator.

% End of Chapter 2. Chapter 3 follows in chapters/chapter3_implementation.tex.
